\chapter{Физический бимодуль заряженных одноформ на пакете \texorpdfstring{$H_{15}$}{H15}}

Предыдущая глава сохранила пятнадцатимерный наблюдаемый пакет и отделила
нейтринный оператор высшей степени. Теперь выполняется тест, ради которого
был переоткрыт путь связности: вместо нуль-форменных эндоморфизмов строится
настоящий модуль одноформ со значениями в трёх заряженных рёбрах Юкавы.

\section{Антикруговая сверка}

Новая постановка отличается от уже закрытых ветвей в трёх отношениях.
Обычная внутренняя флуктуация семейной алгебры либо не видела требуемую
карту, либо возвращала всю $M_3(\mathbb C)$. Поздняя классификация
физического угла изучала $\operatorname{End}_B(E)$, то есть нуль-формы.
Здесь же рассматривается
\begin{equation}
 \operatorname{Hom}_B\left(E,E\otimes_B\Omega_B^1\right)
\end{equation}
после фиксации точного графа $u,d,e$ и ранг-один семейного проектора.

Это различение соответствует общей теории: конечнопорождённый проективный
модуль допускает грассманову связность, а множество всех связностей является
аффинным пространством над модульными отображениями со значениями в
одноформах. В спектральной геометрии Стандартной модели конечный оператор
Дирака задаёт допустимые юкавские блоки, но численные матрицы Юкавы остаются
его входными данными; см. \path{arXiv:hep-th/0608226},
\path{arXiv:1906.09583}, \path{arXiv:1703.05279} и
\path{arXiv:2403.13735}.

\section{Пространство трёх физических рёбер}

Пусть $\Lambda_{\rm ch}$ обозначает пространство кратностей трёх
калибровочно допустимых интертвинеров:
\begin{align}
 \lambda_u&\in\operatorname{Hom}_G(Q_L\otimes\widetilde H,u_R),\\
 \lambda_d&\in\operatorname{Hom}_G(Q_L\otimes H,d_R),\\
 \lambda_e&\in\operatorname{Hom}_G(L_L\otimes H,e_R).
\end{align}
Каждое пространство одномерно по лемме Шура. Однако три рёбра попарно
неэквивалентны как бимодули: у них различаются правые концы, гиперзарядный
канал или цветовое представление. Поэтому
\begin{equation}
 \Lambda_{\rm ch}\simeq\mathbb C^3,
 \qquad
 \operatorname{End}_{\mathcal A_{\rm SM}-\mathcal A_{\rm SM}}
 (\Lambda_{\rm ch})\simeq\mathbb C^3.
 \label{eq:v5-h15-edge-commutant}
\end{equation}
В машинном аудите матрица размерностей интертвинеров между рёбрами равна
$I_3$: внедиагональные пространства морфизмов нулевые.

\section{Правильно типизированный модуль одноформ}

Ранг-один геометрия уже дала семейное пространство представленных одноформ
\begin{equation}
 \mathcal E_\rho=\rho M_3Q\oplus QM_3\rho,
 \qquad \dim_{\mathbb C}\mathcal E_\rho=4.
\end{equation}
Следовательно, пространство кратностей физических заряженных одноформ
имеет вид
\begin{equation}
 \mathcal Y_\rho=\mathcal E_\rho\otimes\Lambda_{\rm ch},
 \qquad \dim_{\mathbb C}\mathcal Y_\rho=4\cdot3=12.
 \label{eq:v5-h15-physical-oneforms}
\end{equation}
Здесь не пересчитываются две компоненты самого хиггсовского дублета:
$\Lambda_{\rm ch}$ является пространством кратностей эквивариантных
рёбер, а не полным пространством значений поля.

Алгебра Стандартной модели конечномерна и полупроста, поэтому её конечные
модули проективны. Прямая сумма трёх простых рёбер также проективна, а
тензорное умножение на уже построенный конечный семейный модуль не нарушает
это свойство. Таким образом, препятствия существованию грассмановой
связности здесь нет.

\section{Что фиксирует грассманова связность}

Ранг-один проектор задаёт эталонную семейную связность
\begin{equation}
 \nabla_\rho^{\rm Gr}=\rho\,d
\end{equation}
на соответствующем проективном модуле. На трёх рёбрах она действует как
\begin{equation}
 \nabla_{\rm ch}^{\rm Gr}=\nabla_\rho^{\rm Gr}\otimes I_{\Lambda_{\rm ch}}.
\end{equation}
Семейная часть связности получает каноническое начало отсчёта и не требует
выбора произвольной матрицы из $M_3(\mathbb C)$.

Но тождественный множитель на $\Lambda_{\rm ch}$ не выбирает относительные
коэффициенты трёх рёбер. Разности совместимых связностей содержат
\begin{equation}
 \alpha=\operatorname{diag}(\alpha_u,\alpha_d,\alpha_e)\in\mathbb C^3.
\end{equation}
Эрмитовость превращает их в три вещественных направления соответствующего
самосопряжённого или косоэрмитова среза. Удаление общего направления
$(1,1,1)$ оставляет двумерное относительное пространство
\begin{equation}
 \mathbb R^3/\mathbb R(1,1,1)\simeq\mathbb R^2.
 \label{eq:v5-h15-relative-connection-plane}
\end{equation}
Условия KO6 спаривают каждое ребро с его сопряжённым, но не отождествляют
$u$, $d$ и $e$, поскольку их бимодульные метки различны.

\begin{theorem}[Одноформенный проход и остаточная двумерная свобода]
На замороженном пакете $H_{15}$ существует проективный физический модуль
заряженных одноформ $\mathcal Y_\rho$ размерности двенадцать в пространстве
кратностей. Грассманова связность канонически фиксирует его семейную часть.
Однако коммутант трёх попарно неэквивалентных рёбер равен $\mathbb C^3$;
после удаления общего направления остаются две вещественные относительные
свободы. Поэтому исправление типа связности само по себе не выводит
единственный юкавский оператор.
\end{theorem}

\begin{nogocriterion}[Запрет отождествления неэквивалентных рёбер]
Нельзя положить $\alpha_u=\alpha_d=\alpha_e$ только на основании общего
хиггсовского дублета или грассмановой семейной связности. Такое равенство
является дополнительной межсекторной симметрией. В текущей алгебре она
отсутствует, а её введение изменило бы наблюдаемое представление.
\end{nogocriterion}

\section{Последняя дешёвая проверка этой щели}

Новая ветвь не должна возвращаться к подбору потенциала на оставшейся
плоскости. Единственный ещё не выполненный внутренний тест --- вычислить
дифференциальное исчисление второй степени и спектральное кручение на прямой
сумме трёх рёбер. Работа \path{arXiv:2511.08159} показывает, что изменение
второй степени способно восстанавливать информацию, невидимую обычным
одноформам. Если соответствующий функционал остаётся блочно диагональным по
$u,d,e$, две относительные свободы сохраняются и одноформенный путь Мориты
следует закрыть без нового потенциала.

\begin{protocol}
\begin{enumerate}
 \item правильно типизированный модуль одноформ: \textbf{построен};
 \item проективность: \textbf{получена из полупростоты конечной алгебры};
 \item ровно три заряженных типа рёбер: \textbf{получены};
 \item грассманова семейная связность: \textbf{получена};
 \item смешивание неэквивалентных рёбер:
 \textbf{запрещено бимодульными метками};
 \item остаточная относительная размерность:
 \textbf{две вещественные свободы};
 \item единственный юкавский оператор: \textbf{не выведен};
 \item физическое замыкание: \textbf{не достигнуто}.
\end{enumerate}
\end{protocol}

Машинный сертификат сохранён в
\path{s2t_v5_h15_physical_oneform_bimodule_gate.py} и
\path{s2t_v5_h15_physical_oneform_bimodule_gate_results.json}.